% GENERATED FILE. Edit canonical lessons, then run npm run generate:books.
% canonical-source-hash: fnv1a64:96b085383fed2f1e
% canonical-lessons: UR-C386-injan, UR-C386-dabba, UR-C386-lain, UR-C386-siti, UR-C386-bajna

\chapter{Engines and Carriages}
\label{ch:ur-engines-and-carriages}
\hlchaptermodality{386}

\begin{chapteropening}
\textbf{By the end of this chapter:} \emph{I can talk about the engine, the carriages and the whistle of a train.}

Complete the last lesson of chapter 386: I can talk about the engine, the carriages and the whistle of a train.
\end{chapteropening}

\section[\texorpdfstring{\ur{انجن} (injan)}{انجن (injan)}]{\ur{انجن} (injan) — an engine}

\label{lesson:UR-C386-injan}



\begin{quote}
Before the new one: say the Urdu for a fruit seller, then the Urdu for a greengrocer.
\end{quote}

\subsection*{You'll want to know: \ur{انجن}}
\textbf{\ur{انجن}} — \emph{injan} — \textquotedblleft{}an engine\textquotedblright{}.

\begin{quote}
\ur{ریل} \ur{گاڑی} \ur{کا} \ur{انجن} \ur{کالا} \ur{اور} \ur{بڑا} \ur{ہے۔}
\end{quote}

\begin{quote}
\emph{rel gāṛī kā injan kālā aur baṛā hai.}
\end{quote}

\begin{quote}
The train's engine is big and black.
\end{quote}

Masculine.

\begin{cousinweb}
\textbf{\ur{انجن}} (\emph{injan}) is borrowed from English \emph{engine}; at a station it means the locomotive at the front of the train.
\end{cousinweb}

\begin{grammarlens}[title={What you've built}]
One more word for getting around.
\end{grammarlens}

\subsection*{Your turn}
\begin{itemize}
\raggedright
  \item \emph{Say it:} \emph{injan}
  \item \emph{Say it:} \emph{injan}, once more
  \item \emph{Say it:} \emph{rel gāṛī kā injan kālā aur baṛā hai.} — The train's engine is big and black.
  \item \emph{Recall:} say the Urdu for a fruit seller, then the Urdu for a greengrocer, then say \emph{injan} again
\end{itemize}

\begin{tcolorbox}[breakable,colback=teal!4,colframe=teal!35!black,title={Before you move on}]
What does \ur{انجن} mean? (\textquotedblleft{}An engine\textquotedblright{}.) Say it once more.
\end{tcolorbox}
\section[\texorpdfstring{\ur{ڈبہ} (ḍabbā)}{ڈبہ (ḍabbā)}]{\ur{ڈبہ} (ḍabbā) — a compartment (of a train); a box}

\label{lesson:UR-C386-dabba}



\begin{quote}
Before the new one: say the Urdu for a greengrocer, then the Urdu for an engine.
\end{quote}

\subsection*{You'll want to know: \ur{ڈبہ}}
\textbf{\ur{ڈبہ}} — \emph{ḍabbā} — \textquotedblleft{}a compartment (of a train); a box\textquotedblright{}.

\begin{quote}
\ur{ہمارا} \ur{ڈبہ} \ur{انجن} \ur{کے} \ur{پیچھے} \ur{ہے۔}
\end{quote}

\begin{quote}
\emph{hamārā ḍabbā injan ke pīchhe hai.}
\end{quote}

\begin{quote}
Our carriage is right behind the engine.
\end{quote}

Masculine.

\begin{cousinweb}
\textbf{\ur{ڈبہ}} (\emph{ḍabbā}) comes from the word CDIAL files as \emph{ḍabba}, \textquotedblleft{}a box\textquotedblright{}, of non-Aryan, perhaps Dravidian, origin (CDIAL 5528); Urdu also calls a railway carriage a box.
\end{cousinweb}

\begin{grammarlens}[title={What you've built}]
One more word for getting around.
\end{grammarlens}

\subsection*{Your turn}
\begin{itemize}
\raggedright
  \item \emph{Say it:} \emph{ḍabbā}
  \item \emph{Say it:} \emph{ḍabbā}, once more
  \item \emph{Say it:} \emph{hamārā ḍabbā injan ke pīchhe hai.} — Our carriage is right behind the engine.
  \item \emph{Recall:} say the Urdu for a greengrocer, then the Urdu for an engine, then say \emph{ḍabbā} again
\end{itemize}

\begin{tcolorbox}[breakable,colback=teal!4,colframe=teal!35!black,title={Before you move on}]
What does \ur{ڈبہ} mean? (\textquotedblleft{}A compartment (of a train); a box\textquotedblright{}.) Say it once more.
\end{tcolorbox}
\section[\texorpdfstring{\ur{لائن} (lāin)}{لائن (lāin)}]{\ur{لائن} (lāin) — a railway line; a line}

\label{lesson:UR-C386-lain}



\begin{quote}
Before the new one: say the Urdu for an engine, then the Urdu for a compartment (of a train).
\end{quote}

\subsection*{You'll want to know: \ur{لائن}}
\textbf{\ur{لائن}} — \emph{lāin} — \textquotedblleft{}a railway line; a line\textquotedblright{}.

\begin{quote}
\ur{یہ} \ur{لائن} \ur{لاہور} \ur{تک} \ur{جاتی} \ur{ہے۔}
\end{quote}

\begin{quote}
\emph{yih lāin lāhaur tak jātī hai.}
\end{quote}

\begin{quote}
This line runs as far as Lahore.
\end{quote}

Feminine.

\begin{cousinweb}
\textbf{\ur{لائن}} (\emph{lāin}) is borrowed from English \emph{line}; Urdu uses it for a railway line and for a telephone line.
\end{cousinweb}

\begin{grammarlens}[title={What you've built}]
One more word for getting around.
\end{grammarlens}

\subsection*{Your turn}
\begin{itemize}
\raggedright
  \item \emph{Say it:} \emph{lāin}
  \item \emph{Say it:} \emph{lāin}, once more
  \item \emph{Say it:} \emph{yih lāin lāhaur tak jātī hai.} — This line runs as far as Lahore.
  \item \emph{Recall:} say the Urdu for an engine, then the Urdu for a compartment (of a train), then say \emph{lāin} again
\end{itemize}

\begin{tcolorbox}[breakable,colback=teal!4,colframe=teal!35!black,title={Before you move on}]
What does \ur{لائن} mean? (\textquotedblleft{}A railway line; a line\textquotedblright{}.) Say it once more.
\end{tcolorbox}
\section[\texorpdfstring{\ur{سیٹی} (sīṭī)}{سیٹی (sīṭī)}]{\ur{سیٹی} (sīṭī) — a whistle}

\label{lesson:UR-C386-siti}



\begin{quote}
Before the new one: say the Urdu for a compartment (of a train), then the Urdu for a railway line.
\end{quote}

\subsection*{You'll want to know: \ur{سیٹی}}
\textbf{\ur{سیٹی}} — \emph{sīṭī} — \textquotedblleft{}a whistle\textquotedblright{}.

\begin{quote}
\ur{میں} \ur{نے} \ur{صبح} \ur{انجن} \ur{کی} \ur{سیٹی} \ur{سنی۔}
\end{quote}

\begin{quote}
\emph{maiṅ ne subah injan kī sīṭī sunī.}
\end{quote}

\begin{quote}
In the morning I heard the engine's whistle.
\end{quote}

Feminine.

\begin{cousinweb}
\textbf{\ur{سیٹی}} (\emph{sīṭī}) is inherited from \emph{sīṭṭa}, \textquotedblleft{}a whistle\textquotedblright{}, a form CDIAL reconstructs and treats as an imitation of the sound (CDIAL 13427).
\end{cousinweb}

\begin{grammarlens}[title={What you've built}]
One more word for getting around.
\end{grammarlens}

\subsection*{Your turn}
\begin{itemize}
\raggedright
  \item \emph{Say it:} \emph{sīṭī}
  \item \emph{Say it:} \emph{sīṭī}, once more
  \item \emph{Say it:} \emph{maiṅ ne subah injan kī sīṭī sunī.} — In the morning I heard the engine's whistle.
  \item \emph{Recall:} say the Urdu for a compartment (of a train), then the Urdu for a railway line, then say \emph{sīṭī} again
\end{itemize}

\begin{tcolorbox}[breakable,colback=teal!4,colframe=teal!35!black,title={Before you move on}]
What does \ur{سیٹی} mean? (\textquotedblleft{}A whistle\textquotedblright{}.) Say it once more.
\end{tcolorbox}
\section[\texorpdfstring{\ur{بجنا} (bajnā)}{بجنا (bajnā)}]{\ur{بجنا} (bajnā) — to sound; to ring}

\label{lesson:UR-C386-bajna}



\begin{quote}
Before the new one: say the Urdu for a railway line, then the Urdu for a whistle.
\end{quote}

\subsection*{You'll want to know: \ur{بجنا}}
\textbf{\ur{بجنا}} — \emph{bajnā} — \textquotedblleft{}to sound; to ring\textquotedblright{}.

\begin{quote}
\ur{اسٹیشن} \ur{پر} \ur{سیٹی} \ur{بجی} \ur{اور} \ur{گاڑی} \ur{چلی۔}
\end{quote}

\begin{quote}
\emph{isṭeshan par sīṭī bajī aur gāṛī chalī.}
\end{quote}

\begin{quote}
The whistle sounded at the station and the train moved off.
\end{quote}

Past (it …): \ur{وہ} \ur{سیٹی} \ur{بجی}, \emph{voh sīṭī bajī}.

\begin{cousinweb}
\textbf{\ur{بجنا}} (\emph{bajnā}) is inherited: CDIAL lists the Hindi cousin \emph{bājnā}, \textquotedblleft{}to sound\textquotedblright{}, under Sanskrit \emph{vādyate}, \textquotedblleft{}is made to sound\textquotedblright{} (CDIAL 11513). \textbf{\ur{بجانا}}, \textquotedblleft{}to play an instrument\textquotedblright{}, is the form for making a thing sound, and \textbf{\ur{بجے}}, \textquotedblleft{}o'clock\textquotedblright{}, is the hour that has struck.
\end{cousinweb}

\begin{grammarlens}[title={What you've built}]
One more word for getting around.
\end{grammarlens}

\subsection*{Your turn}
\begin{itemize}
\raggedright
  \item \emph{Say it:} \emph{bajnā}
  \item \emph{Say it:} \emph{bajnā}, once more
  \item \emph{Say it:} \emph{isṭeshan par sīṭī bajī aur gāṛī chalī.} — The whistle sounded at the station and the train moved off.
  \item \emph{Recall:} say the Urdu for a railway line, then the Urdu for a whistle, then say \emph{bajnā} again
\end{itemize}

\begin{tcolorbox}[breakable,colback=teal!4,colframe=teal!35!black,title={Before you move on}]
What does \ur{بجنا} mean? (\textquotedblleft{}To sound; to ring\textquotedblright{}.) Say it once more.
\end{tcolorbox}
